%!TEX root = main.tex

\subsection{Formalization of Transformations}

%Equipped with this definition of
With the formal semantics of {\hifirrtl} programs, we %are able to
shall formalize the lowering transformations and prove %that they preserve the
various semantics preservation properties in Section~\ref{sec-hi-firrtl-transform}. %, i.e.,  %of {\hifirrtl} programs, that is, prove that
%the set of module graphs of {\hifirrtl} programs are unchanged during the lowering transformations.

In this section, we show how the lowering transformations of HiFIRRTL, which were described informally in Section~\ref{subsec-lowering-transformations}, are formalized in Coq. Specifically, we consider the formalization of the six lowering transformations, namely, %``InferTypes'',
``InferWidths'', %``InferResets'',
``ExpandConnects'', and ``ExpandWhens''. %, and ``LowerTypes''.




For simplicity, let us focus on a HiFIRRTL circuit $A$ that contains no \prettyll{inst of} statements, that is, circuit $A$ contains exactly one top-level module $A$.

%The formalization utilizes a structure called \emph{circuit diagrams} defined below.
%
%\begin{definition}[Circuit diagrams]\label{def-cd}
%A \emph{circuit diagram} $D$ is a triple $(ss, ce, cm)$, where $ss$ is a sequence of statements, $ce$ is a component environment that maps each component to its type, and $cm$ is a connect mapping that maps each component to an expression that is connected to it.
%\end{definition}
%
%Note that the types in the range of the component environment $ce$ in Definition~\ref{def-cd} are not necessarily ground types, that is, they are aggregate types in general, since HiFIRRTL may contain high-level constructs, e.g. vectors and bundles. Intuitively, $ss$ is the sequence of statements in a HiFIRRTL program, while $ce$ and $cm$ extract some necessary information from
%the sequence $ss$ of statements to facilitate the formalization of the lowering transformations. Specifically,
%the component environment $ce$ extracts the type information from the sequence $ss$ of statements and
%the connect mapping $cm$ extracts for each component the expression that is connected to the component according to the last-connect semantics.


\tl{to be rewritten}

As we will see later on, the sequence $ss$ of statements will be updated by the ``ExpandConnects'', ``ExpandWhens'', and ``LowerTypes'' passes,
and the component environment $ce$ will be modified by the ``InferTypes'', ``InferWidths'', ``InferResets'',  and ``LowerTypes'' passes,
 while the component environment $cm$ will only be changed by the ``ExpandWhens'' pass, since it is the only pass that deals with the last-connect semantics.

Before applying the lowering transformations, the module $A$ will be first transformed into a circuit diagram $D_A = (ss_0, ce_0, cm_0)$, where $ss_0$ is the sequence of statements in the module  $A$, $ce_0$ is the component environment obtained from the component declarations in the module $A$, possibly with the types of some components missing or partial, and $cm_0$ is the trivial mapping that maps each component to $\bot$. (Intuitively, this means that none of the components have been connected to so far.) The missing or partial types in the component environment $ce_0$ will be completed by the lowering transformations ``InferTypes'' and ``InferWidths''.



%In the sequel, we utilize the concept of circuit diagrams to illustrate how the six lowering transformations are formalized in Coq.

%A HiFIRRTL module $M$ will be first transformed into a circuit diagram $D_M = (ss, ce, cm)$, where $ss$ is the sequence of statements in $M$, $ce$ is the component environment obtained from the variable declarations in $M$, possibly with the types of some variables missing or partial, and $cm$ is the trivial mapping that maps each component to $\bot$ (Intuitively, this means that the connects of the components have not be determined yet). The missing or partial  types in $ce$ will be completed by the compilation passes during the transformation of $M$ into a LoFIRRTL module. Moreover, the connect mapping $cm$ will be updated by the ExpandWhens compilation pass. Furthermore, the statement sequence $ss$ will be updated by ExpandConnects, ExpandWhens, LowerTypes, and RemoveReset, while they keep unchanged during the other compilation passes.

%A description of the formalization of InferTypes, InferWidths, InferResets, ExpandConnects, ExpandWhens, LowerTypes, and RemoveResets.

\subsubsection{Formalization of InferTypes}
%Initially, the component environment $ce$ of a circuit diagram $D_M$
%is empty.  The Infer Types pass is then applied to fill $ce$ with
%all the components that are defined in the statement sequence $ss$,
%along with their component kinds and data types.
We formalize the ``InferTypes'' pass by defining a Coq function {\tt infer\_type}, which iterates over the statements in the sequence $ss$ and updates the component environmen $ce$.
%For each component defined in $ss$, the {\tt infer\_type} function will update $ce$
%accordingly.
For instance, the following rule specifies that when the function {\tt infer\_type} is applied to a wire statement \prettyll{wire} $x: t$, it updates the component environmen $ce$ by assigning $(t, Wire)$ to
the wire $x$ (where $Wire$ denotes the kind of the component $x$), resulting into the component environment $ce'$,
%if there is a \prettyll{wire} statement,
%it is added to component environment by
  \begin{gather*}
    \infer[]
    {\seqq{ce \vdash {\tt infer\_type} (\mbox{\prettyll{wire} } x: t) \rightsquigarrow ce'}}
    {\seqq{ce'= ce \uplus \{x : (t, Wire)\}}}.
    \end{gather*}
Moreover, the function {\tt infer\_type} behaves similarly on \prettyll{reg}, \prettyll{input} and \prettyll{output} port statements.
Finally, on a node statement \prettyll{node } $x = e$, {\tt infer\_type} first calculates the type of the expression $e$, denoted by $t_e$, then assigns $(t_e, Node)$ to
the node $x$, resulting into the component environment $ce'$,
%
%If it is a node-declaration statement, an entry for the node is added
%to $ce$ with its type calculated by the expression $e$
%assigned to the node.
\begin{gather*}
    \inferii[]
    {\seqq{ce \vdash {\tt infer\_type} (\mbox{\prettyll{node} } v = e) \rightsquigarrow ce'}}
    {\seqq{ce \vdash t_e={\tt type\_of\_expr}(e)}} {\seqq{ce' = ce \uplus \{v : (t_e, Node)\}}}.
    \end{gather*}
Note that the function {\tt infer\_type} will not update the component environment $ce$ when it is applied to connect statements.
Nevertheless, in this case, as specified in the following two rules, it checks whether the two types of both sides of the connect statement are equivalent. By type equivalence, we mean that the structures of the types are the same, but their bit widths may be missing or different. For instance, the types \prettyll{UInt}, \prettyll{UInt<1>}, and \prettyll{UInt<2>} are seen as equivalent.
If the two types are equivalent, then the component environment $ce$ is kept, otherwise, an error (indicated by $\bot$) will be issued.
\begin{gather*}
    \inferiii[]
    {\seqq{ce \vdash {\tt infer\_type} (v<=e) \rightsquigarrow ce}}
    {\seqq{ce(v) = (t_v,-)}}
    {\seqq{ce \vdash t_e={\tt type\_of\_expr}(e)}} {\seqq{{\tt equiv}(t_v,t_e)}}
    \end{gather*}
%If the type equivalence is violated, an error will be raised.
\begin{gather*}
    \inferiii[]
    {\seqq{ce \vdash {\tt infer\_type} (v<=e) \rightsquigarrow \bot}}
    {\seqq{ce(v) = (t_v, -)}}
    {\seqq{ce \vdash t_e={\tt type\_of\_expr}(e)}} {\seqq{\neg{\tt equiv}(t_v,t_e)}}
    \end{gather*}
After the function {\tt infer\_type} has been applied to all the statements in the sequence $ss$, if the bit width of a component is still missing, then the zero width will be added for the component in
the component environment $ce$.
%After applying the Infer Types pass, types with an
%implicitly defined width are assigned with 0 in $ce$.

\subsubsection{Formalization of InferWidths}
%% %%%!TEX root = main.tex

%\subsubsection{Formalization of InferWidths}
%By Keyin
We formalize the ``InferWidths'' pass by defining a Coq function ${\sf inferWidthsM}:{\sf M} \rightarrow {\sf M}$, which contains two process. The first process ${\sf inferWidthsStage1 : M \rightarrow tmap}$ calculates the topological order of dependency graph of $M$ and iterates on the topological sorted order to updates a $tmap$, so that the bit widths of all the components with undefined bit widths can be determined. Then the second process ${\sf inferWidthsStage2 : M \times tmap \rightarrow M}$ produce a new module with inferred bit widths and explicit widths instead of implicit widths in the original module.

The rationale behind the {\sf inferWidthM} function is stated as follows: The function  {\sf inferWidthsStage1} always chooses the maximum width of the expressions that are connected to the component $v$. 

The behavior of the function {\sf inferwidthsM} on a module contains the folowing steps :
\begin{enumerate}
    \item Record the original $tmap$ with aggerate types and the connection map $var2exprs{\sf : range(}id{\sf )} \rightarrow \exps$ with only ground type expressions according to ports sequecnse $pp$ and statement sequence $ss$ in the module $(pp,ss)$. 

The Coq function ${\sf {\sf portsStmtsTmap'} : }(pp,ss){\sf \rightarrow}(tmap,var2exprs)$ shown in ? and ?, given ports list and statements list in the module, generates two corresponding maps, mapping from aggregate level components to types and ground type level components to expressions. Note that aggregate connections are not expanded yet so that flip field should be considered when generating the connection map. This is shown in the function ${\sf listGRef :\ }id \times \ftypes \times {\sf Boolean} \rightarrow {\sf seq\ (} id \times \gtypes \times {\sf Boolean} {\sf )}$, where given a identifier, its {\sf FType} and its flip case, generates a list of triples, including identifier of a {\sf GType}, the {\sf GType} and the flip case of the {\sf GType}. The Coq function ${\sf typeOfRef' : } id \times tmap \rightarrow {\sf FType} \times Boolean$ shown in ?  given identifier and $tmap$, generates a pair with the first element is {\sf FType} and the second element is a Boolean value shows whether the component is flipped or not.

For a statement that declares a component $v$ of \datatype \ $\{{\sf flip}\ a : \uint, b : \sint\}$,
the {\sf {\sf portsStmtsTmap'}} function will update $tmap$ by
 $$ {tmap':= tmap \uplus \{v : \{{\sf flip} \  a : \uint, b : \sint\}\}} $$
\\For a statement that connects $v$ to another component $w$ of \datatype \ $\{{\sf flip}\ a : \uint, b : \sint\}$,
the {\sf {\sf portsStmtsTmap'}} function will update $var2exprs$ by
 $$ {var2exprs':= var2exprs \uplus \{v.a : \{w.a\} \} \uplus \{v.b : \{w.b\} \}} $$

    \item Draw the dependency graph $G$ according to the connection map. The Coq function ${\sf drawG :} var2expr \rightarrow G$, given elements in the connection map $var2exprs$, generates the dependency graph by adding edges from every component in a expression to the target component. To make the calculation simpler, only components with uninferred implicit width are involved in the dependency graph,
    
    
    \item Generate the topological order of all the uninferred ground type components. Given the vertices set and the graph, the correctness of sorting algorithm ${\sf topoSort : } G \rightarrow {\sf seq} id $ is certified. Because the module components depends on each other through connections and operations, implicit widths should be inferred sequentially in the order of topological order,

    \item Compare the widths of each expression connected to the same component, and chooses the larger width. Since the last connect semantics is not considered, each connection to the component provides a potential width. A function named ${\sf maxFtList : GType \times \gtypes \rightarrow GType}$ compares the potential widths and outputs the maximum one. Finally, we update the corresponding widths of $v$ in $tmap$.
    
the Coq function {\sf inferWidthsFun} calls {\sf {\sf inferWidthsFun}} on the list of topological soreted order, where the widths are inferred sequentially according to the dependencies. 
\end{enumerate}

\begin{figure}[htbp]
  \begin{gather*}
     \inferii[ ]
    {\seqq{{\sf portsStmtsTmap'}(pp, ss) = (tmap', var2exprs)}}
    {\seqq{{\sf portsTmap'}(pp) = pmap'}}
    {\seqq{{\sf stmtsTmap'}(pmap', ss) = (tmap', var2exprs)}}
   \quad
   \\[1ex]
    \infer[{\sf portsTmap'}: {\sf input}]
    {\seqq{{\sf portsTmap'}(pp) = pmap'}}
    {
    \begin{array}{c}
    p p= {\sf input}\ x: t\ pp' \wedge {\sf portsTmap'}(pp') = pmap \wedge pmap(id(x)) = \svnone \\
    pmap' = \fadd((id(x), (t, \source)), pmap)
    \end{array}
    }
    \quad
    \\[1ex]
    \infer[{\sf portsTmap'}: {\sf output}]
    {\seqq{{\sf portsTmap'}(pp) = pmap'}}
    {
    \begin{array}{c}
    p p= {\sf output}\ x: t\ pp' \wedge {\sf portsTmap'}(pp') = pmap \wedge pmap(id(x)) = \svnone \\
    pmap' = \fadd((id(x), (t, \sink)), pmap)
    \end{array}
    }
    \quad
    \\[1ex]
    \inferii[{\sf stmtsTmap'}: \nil]
    {\seqq{{\sf stmtsTmap'}((tmap, var2exprs), ss) = (tmap', var2exprs')}}
    {\seqq{s s = \nil}}
    {\seqq{tmap' = tmap \wedge var2exprs' = var2exprs}}
    \quad
    \\[1ex]
    \infer[{\sf stmtsTmap'}: $\neq$ \nil]
    {\seqq{{\sf stmtsTmap'}((tmap, var2exprs), ss) = (tmap', var2exprs')}}
    {
    \begin{array}{c}
      ss = s\ ss' \\
      (tmap'', var2exprs'') = stmtTmap((tmap, var2exprs), s) \\
      (tmap', var2exprs') = stmtsTmap((tmap'', var2exprs''), ss')
    \end{array}
    }
  \end{gather*}
%  \vspace{-4mm}
  \caption{${\sf portsStmtsTmap'}$, ${\sf portsTmap'}$, and ${\sf stmtsTmap'}$}
  \label{fig:ports_stmts_tmap'}
\end{figure}

\begin{figure}[htbp]
  \begin{gather*}
   \infer[find-some]
    {\seqq{{\sf addExprConnect}(lhsl, rhsl, var2exprs) = var2exprs'}}
    {
    \begin{array}{c}
    lhsl = (ref1, \_, \_) \ ltl \wedge rhsl = ehd \ etl \\
    var2exprs(ref1) = el \wedge var2exprs'' = \fadd(ref1, ehd :: el, var2exprs) \\
    var2exprs' = {\sf addExprConnect}(ltl, etl, var2exprs'')
    \end{array}
    }
    \quad
    \\[1ex]
    \infer[find-none]
    {\seqq{{\sf addExprConnect}(lhsl, rhsl, var2exprs) = var2exprs'}}
    {
    \begin{array}{c}
    lhsl = (ref1, \_, \_) \ ltl \wedge rhsl = ehd \ etl \\
    var2exprs(ref1) = \svnone \wedge var2exprs'' = \fadd(ref1, [::ehd], var2exprs) \\
    var2exprs' = {\sf addExprConnect}(ltl, etl, var2exprs'')
    \end{array}
    }
     \quad
    \\[1ex]
    \infer[flip, find-none]
    {\seqq{{\sf addExprConnectNonPassive}(lhsl, rhsl, var2exprs) = var2exprs'}}
    {
    \begin{array}{c}
    lhsl = (ref1, \_, true) \ ltl \wedge rhsl = (ref2, \_, \_) \ rtl \\
    var2exprs(ref2) = \svnone \wedge var2exprs'' = \fadd(ref2, [::(Eref \ ref1)], var2exprs) \\
    var2exprs' = {\sf addExprConnectNonPassive}(ltl, rtl, var2exprs'')
    \end{array}
    }
     \quad
    \\[1ex]
    \infer[non-flip, find-none]
    {\seqq{{\sf addExprConnectNonPassive}(lhsl, rhsl, var2exprs) = var2exprs'}}
    {
    \begin{array}{c}
    lhsl = (ref1, \_, false) \ ltl \wedge rhsl = (ref2, \_, \_) \ rtl \\
    var2exprs(ref1) = \svnone \wedge var2exprs'' = \fadd(ref1, [::(Eref \ ref2)], var2exprs) \\
    var2exprs' = {\sf addExprConnectNonPassive}(ltl, rtl, var2exprs'')
    \end{array}
    }
    \quad
    \\[1ex]
    \infer[flip, find-some]
    {\seqq{{\sf addExprConnectNonPassive}(lhsl, rhsl, var2exprs) = var2exprs'}}
    {
    \begin{array}{c}
    lhsl = (ref1, \_, true) \ ltl \wedge rhsl = (ref2, \_, \_) \ rtl \\
    var2exprs(ref2) = el \wedge var2exprs'' = \fadd(ref2, (Eref \ ref1) :: el, var2exprs) \\
    var2exprs' = {\sf addExprConnectNonPassive}(ltl, rtl, var2exprs'')
    \end{array}
    }
    \quad
    \\[1ex]
    \infer[non-flip, find-some]
    {\seqq{{\sf addExprConnectNonPassive}(lhsl, rhsl, var2exprs) = var2exprs'}}
    {
    \begin{array}{c}
    lhsl = (ref1, \_, false) \ ltl \wedge rhsl = (ref2, \_, \_) \ rtl \\
    var2exprs(ref1) = el \wedge var2exprs'' = \fadd(ref1, (Eref \ ref2) :: el, var2exprs) \\
    var2exprs' = {\sf addExprConnectNonPassive}(ltl, rtl, var2exprs'')
    \end{array}
    }
  \end{gather*}
\caption{${\sf addExprConnectNonPassive}$ and ${\sf addExprConnect}$}
  \label{tab:add_expr_connect}
\end{figure}


%%%%% where the bug begins
\begin{figure}[htbp]
   \begin{gather*}
    \infer[{\sf stmtTmap'}: skip]
    {\seqq{{\sf stmtTmap'}((tmap, var2exprs), s)=(tmap, var2exprs)}}
    {\seqq{s = skip}}
   \quad
   \\[1ex]
    \infer[{\sf stmtTmap'}: connect-ref]
    {\seqq{{\sf stmtTmap'}((tmap, var2exprs), s) = (tmap, var2exprs')}}
    {
    \begin{array}{c}
      s = r <= ref \\
      {\sf typeOfRef'}(r, tmap) = (t\_tgt, flip\_tgt)  \\
      {\sf typeOfRef'}(ref, tmap) = (t\_src, flip\_src) \\ 
      lhsl = {\sf listGRef} (r, t\_tgt, flip\_tgt) \\
      rhsl = {\sf listGRef} (ref, t\_src, flip\_src) \\
      var2exprs' = {\sf addExprConnectNonPassive} ((lhsl, rhsl), var2exprs) 
     \end{array}
    }
   \quad
   \\[1ex]
    \infer[{\sf stmtTmap'}: connect-expr]
    {\seqq{{\sf stmtTmap'}((tmap, var2exprs), s) = (tmap, var2exprs')}}
    {
    \begin{array}{c}
      s = r <= exp \\
      {\sf typeOfRef'}(r, tmap) = (t\_tgt, \_)  \\
      {\sf typeOfExp}(exp, tmap) = t\_src \\ 
      lhsl = {\sf listGRef} (r, t\_tgt, \lfalse) \\
      rhsl = {\sf listGExp} (exp, t\_src)\\
      var2exprs' = {\sf addExprConnect} ((lhsl, rhsl), var2exprs)
     \end{array}
    }
   \quad
   \\[1ex]
    \infer[{\sf stmtTmap'}: invalid]
    {\seqq{{\sf stmtTmap'}((tmap, var2exprs), s) = (tmap, var2exprs)}}
    {
      \begin{array}{c}
      s = r \mbox{ is invalid} \\
      {\sf typeOfRef'}(r, tmap) \neq \svnone
     \end{array}
    }
   \quad
   \\[1ex]
    \infer[{\sf stmtTmap'}: wire]
    {\seqq{{\sf stmtTmap'}((tmap, var2exprs), s) = (tmap', var2exprs)}}
    {
      \begin{array}{c}
      s = \mbox{wire } x: t \wedge tmap(id(x)) = \svnone  \\
      tmap' = \fadd((id(x), (t, \duplex)), tmap) 
     \end{array}
    }
   \quad
   \\[1ex]
    \infer[{\sf stmtTmap'}: node]
    {\seqq{{\sf stmtTmap'}((tmap, var2exprs), s) = (tmap', var2exprs')}}
    {
     \begin{array}{c}
      s = \mbox{node } x = e \wedge tmap(id(x)) = \svnone  \\
      {\sf typeOfExp}(e, tmap) = t\_src \\
      tmap' = \fadd((id(x), (t\_src, \source)), tmap) \\
      lhsl = {\sf listGRef} (x, t\_src, false) \\
      rhsl = {\sf listGExp} (e, t\_src)\\
      var2exprs' = {\sf addExprConnect} ((lhsl, rhsl),var2exprs )
     \end{array}
    }
\end{gather*}
%  \vspace{-4mm}
  \caption{{\sf stmtTmap'} (part 1)}
  \label{fig:stmts-tmap'}
\end{figure}

\begin{figure}[htbp]
  \small
   \begin{gather*}
    \infer[{\sf stmtTmap'}: register-reset]
    {\seqq{{\sf stmtTmap'}((tmap, var2exprs), s) = (tmap', var2exprs')}}
    {
    \begin{array}{c} 
      s = \mbox{reg } x: t \ clk \mbox{ with: reset =>} (rsig, rval)  \\
      tmap(id(x)) = \svnone \wedge {\sf typeOfExp}(clk, tmap) \neq \svnone  \\
      tmap' = \fadd((id(x), (t, \duplex)), tmap) \\
      lhsl = {\sf listGRef} (x, t, false) \\
      rhsl = {\sf listGExp} (rval, t)\\
      var2exprs' = {\sf addExprConnect} ((lhsl, rhsl), var2exprs )
     \end{array}
    }
   \quad
   \\[1ex]
    \infer[{\sf stmtTmap'}: register-non-reset]
    {\seqq{{\sf stmtTmap'}((tmap, var2exprs), s) = (tmap', var2exprs)}}
    {
    \begin{array}{c}
      s = \mbox{reg } x: t \ clk \\
      tmap(id(x)) = \svnone \wedge {\sf typeOfExp}(clk, tmap) \neq \svnone  \\
      tmap' = \fadd((id(x), (t, \duplex)), tmap) 
     \end{array}
    }
   \quad
   \\[1ex]
   \infer[{\sf stmtTmap'}: when]
    {\seqq{{\sf stmtTmap'}((tmap, var2exprs), s) = (tmap', var2exprs')}}
    {
    \begin{array}{c}
      s = \mbox{when } c: sst \mbox{ else }: ssf  \\
      {\sf typeOfExp}(c, scope) \neq \svnone \\
      {\sf stmtsTmap'}((tmap, var2exprs), sst) = (tmap'', var2exprs'') \\
      {\sf stmtsTmap'}((tmap'', var2exprs''), ssf) = (tmap', var2exprs')
     \end{array}
    }
    \end{gather*}
%  \vspace{-4mm}
  \caption{${\sf stmtTmap'} (part 2)$}
  \label{tab:stmts-tmap-2}
\end{figure}

\begin{figure}[htbp]
  \small
  \begin{gather*}
   \inferii[{\sf typeOfRef'}: id]
    {\seqq{{\sf typeOfRef'}(r, tmap) = (ft, false)}}
    {\seqq{r = x}} {\seqq{tmap(id(x)) = (ft, \_)}}
    \quad
    \\[1ex]
   \inferiii[{\sf typeOfRef'}: vector-1]
    {\seqq{{\sf typeOfRef'}(r, tmap) = (t, flip)}}
    {\seqq{r = r'[n]}} {\seqq{{\sf typeOfRef'}(r', tmap) =  (t[m], flip)}} {\seqq{n < m}}
    \quad
    \\[1ex]
   \inferiii[{\sf typeOfRef'}: vector-2]
    {\seqq{{\sf typeOfRef'}(r, tmap) = \svnone}}
    {\seqq{r = r'[n]}} {\seqq{{\sf typeOfRef'}(r', tmap) = (t[m], flip)}} {\seqq{n \ge m}}
    \quad
    \\[1ex]
   \infer[{\sf typeOfRef'}: field]
    {\seqq{{\sf typeOfRef'}(r, tmap) = {\sf typeOfField'}(f, \myvec{ff'}, flip')}}
    {
    \begin{array}{c}
    r = r'.f \\
    {\sf typeOfRef'}(r', tmap) = ((\btype, \myvec{ff'}), flip')
    \end{array}
    }
    \quad
    \\[1ex]
   \infer[{\sf typeOfField'} : \nil]
    {\seqq{{\sf typeOfField'}(f, \myvec{ff}, flip) = \svnone}}
    {\seqq{\myvec{ff} = \nil}}
    \quad
    \\[1ex]
   \infer[{\sf typeOfField'}: neq]
    {\seqq{{\sf typeOfField'}(f, \myvec{ff}, flip) = {\sf typeOfField'}(f, \myvec{ff'}, flip)}}
    {
    \begin{array}{c}
    \myvec{ff} = (f', flip, t, \myvec{ff'})\wedge f \neq f'
    \end{array}
    }
    \quad
    \\[1ex]
   \infer[{\sf typeOfField'}: eq-1]
    {\seqq{{\sf typeOfField'}(f, \myvec{ff}, flip) = (t, flip)}}
    {
    \begin{array}{c}
    \myvec{ff} = (f', flip, t, \myvec{ff'}) \wedge f = f' \\
    flipped = false
    \end{array}
    }
    \quad
    \\[1ex]
   \infer[{\sf typeOfField'}: eq-2]
    {\seqq{{\sf typeOfField'}(f, \myvec{ff}, flip) = (t, \neg flip)}}
    {
    \begin{array}{c}
    \myvec{ff} = (f', flip, t, \myvec{ff'}) \wedge f = f' \\
    flip = true
    \end{array}
    }
  \end{gather*}
%  \vspace{-4mm}
  \caption{${\sf typeOfRef'}$ and ${\sf typeOfField'}$}
  \label{tab:type-ref'}
\end{figure}


\begin{table}[htbp]
  \small
  \begin{gather*}
    \infer[{\sf drawG}: \nil]
    {\seqq{{\sf drawG}({\sf elements} (var2exprs), tmap) = (\nil, emptyg)}}
    {\seqq{{\sf elements} (var2exprs) = \nil}}
    \quad
   \\[1ex]
    \infer[{\sf drawG}: $\neq$ \nil, explicit]
    {\seqq{{\sf drawG}({\sf elements} (var2exprs), tmap, vertices, g) = {\sf drawG} (tl, tmap, vertices, g)}}
    {
    \begin{array}{c}
      {\sf elements} (var2exprs) = v : \_\ tl \\
      {\sf typeOfRef'}(v, tmap) = (t, \_) \wedge {\sf notImplict} (t) = \ltrue
    \end{array}
    }
    \quad
   \\[1ex]
    \infer[{\sf drawG}: $\neq$ \nil, implicit]
    {\seqq{{\sf drawG}({\sf elements} (var2exprs), tmap, vertices, g) = {\sf drawG} (tl, tmap, vertices'', g'')}}
    {
    \begin{array}{c}
      {\sf elements} (var2exprs) = v : el\ tl \\
      {\sf typeOfRef'}(v, tmap) = (t, \_) \wedge {\sf notImplict} (t) = false \\
      (vertices'', g'') = {\sf drawEl}(v, el, vertices, g) 
    \end{array}
    }
    \quad
   \\[1ex]
    \infer[{\sf drawEl}: \nil]
    {\seqq{{\sf drawEl}(v, el, vertices, g) = (vertices, g)}}
    {\seqq{el = \nil}}
    \quad
   \\[1ex]
    \infer[{\sf drawEl}: $\neq$ \nil]
    {\seqq{{\sf drawEl}(v, el, vertices, g) = {\sf drawEl}(v, tl, vertices', g')}}
    {
    \begin{array}{c}
      el = e\ tl \\
      vl = {\sf expr2varlist} (e) \wedge vertices' = vertices ++ vl \\
      g' = {\sf foldLeft}({\sf fun} (tv, tg) => \fadd(v, tv :: tg(v), tg)) \leftarrow (vl,g)
    \end{array}
    }
    \quad
    \\[1ex]
    \infer[{\sf expr2varlist}: const]
    {\seqq{{\sf expr2varlist}(e) = \nil}}
    {\seqq{e = const}}
   \quad
   \\[1ex]
    \infer[{\sf expr2varlist}: ref]
    {\seqq{{\sf expr2varlist}(e) = [::ref]}}
    {\seqq{e = ref}}
   \quad
   \\[1ex]
    \inferii[{\sf expr2varlist}: cast]
    {\seqq{{\sf expr2varlist}(e) = vl}}
    {\seqq{e = cast(e0)}}
    {\seqq{vl = {\sf expr2varlist}(e0)}}
   \quad
   \\[1ex]
    \inferii[{\sf expr2varlist}: unop]
    {\seqq{{\sf expr2varlist}(e) = vl}}
    {\seqq{e = unop(e0)}}
    {\seqq{vl = {\sf expr2varlist}(e0)}}
   \quad
   \\[1ex]
    \inferii[{\sf expr2varlist}: binop]
    {\seqq{{\sf expr2varlist}(e) = vl}}
    {\seqq{e = binop(e0, e1)}}
    {\seqq{vl = {\sf expr2varlist}(e0) ++ {\sf expr2varlist}(e1)}}
    \quad
   \\[1ex]
    \inferii[{\sf expr2varlist}: mux]
    {\seqq{{\sf expr2varlist}(e) = vl}}
    {\seqq{e = mux(e0, e1, e2)}}
    {\seqq{vl = {\sf expr2varlist}(e0) ++ {\sf expr2varlist}(e1) ++ {\sf expr2varlist}(e2)}}
  \end{gather*}
%  \vspace{-4mm}
  \caption{${\sf drawG}$, ${\sf drawEl}$, and ${\sf expr2varlist}$}
  \label{tab:ports-stmts-tmap}
  \vspace{-6mm}
\end{table}


\begin{table}[htbp]
  \small
  \begin{gather*}
    \infer[{\sf inferWidthsFun}: \nil]
    {\seqq{{\sf inferWidthsFun}(order, var2exprs, tmap) = tmap}}
    {\seqq{order = \nil}}
    \quad
   \\[1ex]
    \infer[{\sf inferWidthsFun}: $\neq$ \nil]
    {\seqq{{\sf inferWidthsFun}(order, var2exprs, tmap) = tmap'}}
    {
    \begin{array}{c}
      order = v\ tl \\
      el = var2exprs(v) \\
      tmap'' = {\sf inferWidthsFun}(v, el, tmap) \\
      tmap' = {\sf inferWidthsFun} (tl, var2exprs, tmap'')
    \end{array}
    }
    \quad
   \\[1ex]
    \infer[{\sf inferWidthsFun}]
    {\seqq{{\sf inferWidthsFun}(v, el, tmap) = tmap'}}
    {
    \begin{array}{c}
      {\sf typeOfRef'}(v, tmap) = (initt, \_) \\
      newt = {\sf maxFtList}(initt, map (({\sf typeOfExp}(tmap)) \leftarrow el)) \\
      tmap' = \fadd(v, newt, tmap)
    \end{array}
    }
    \quad
   \\[1ex]
    \infer[{\sf maxFtList}: \nil]
    {\seqq{{\sf maxFtList}(initt, tel) = initt}}
    {\seqq{tel = \nil}}
    \quad
   \\[1ex]
    \infer[{\sf maxFtList}: $\neq$ \nil]
    {\seqq{{\sf maxFtList}(initt, tel) = {\sf maxFtList}(newt', tl) }}
    {
    \begin{array}{c}
      tel = gt \ tl \\
      newt' = {\sf maxFgtyp}(initt, gt) 
    \end{array}
    }
    \quad
   \\[1ex]
    \infer[{\sf maxFgtyp}: {\sf UIntImp}]
    {\seqq{{\sf maxFgtyp}(initt, te) = ({\sf UIntImp},{\sf max}(w, {\sf sizeOfGtyp}(te)))}}
    {\seqq{initt = ({\sf UIntImp},w)}}
    \quad
   \\[1ex]
    \infer[{\sf maxFgtyp}: {\sf SIntImp}]
    {\seqq{{\sf maxFgtyp}(initt, te) = ({\sf SIntImp},{\sf max}(w, {\sf sizeOfGtyp}(te)))}}
    {\seqq{initt = ({\sf SIntImp},w)}}
  \end{gather*}
 \vspace{-4mm}
  \caption{${\sf inferWidthsFun}$, ${\sf inferWidthsFun}$, ${\sf maxFtList}$ and ${\sf maxFgtyp}$}
  \label{tab:InferWidths-fun}
  \vspace{-6mm}
\end{table}


%%%%%%%%%%%%%%%
%\hide{
%It is allowed to leave the width of a component implicit in HiFIRRTL programs.
%Such unspecified widths are inferred by the Infer Widths pass.
%The general rules are,
%\begin{itemize}
%\item The only allowed width for an explicitly defined width is itself.
%\item An implicitly defined width will be set to the largest width that can
%  be found through connect statements.
%\item The last connection semantics is not considered at this stage,
%  all connect statements are taken into account.
%\end{itemize}
%
%Following the rules, we first introduce an additional map $wm$, which
%maps from variables to types. Then we scan through the statements in $ss$ to
%find components with implicit width and create entries for
%them in $wm$. For a statement that declares a component $v$,
%the {\sf infer\_width} function will update $wm$ by
%%
%  \begin{gather}
%    \inferii[]
%    {\seqq{ce \vdash {\sf infer\_width} (declare\;v, wm) \rightsquigarrow wm'}}
%    {\seqq{ce \vdash {\sf width}(t_v) = 0}} {\seqq{wm':= wm \uplus \{v : t_v\}}}.
%    \end{gather}
%
%The width is inferred according to connect statements.
%Since the last connect semantics is not considered, each
%connection to the element provides a potential width. A function
%named {\sf max\_width} compares the potential widths and outputs
%the maximum one. For a statement which connects an expression $e$
%to a component $v$, the {\sf infer\_width} function will update $wm$ by
%
%After applying {\sf infer\_width} to all the statements in $ss$,
%$wm$ contains the inferred width for every implicitly defined
%component. By merging the information from $wm$ for every entry into $ce$,
%the component environment $ce$ contains the
%result of the inferred width.
%
%Instead of applying {\sf infer\_width} directly on $ce$, the map
%$wm$ contains only type information of the implicit ones that can be
%manipulated more efficiently.
%}
%%%%%%%%%%%%%%%

%% \begin{description}
%% \item[The {\sf input}] of InferWidth takes the statments and the component
%%   environment $ce$ that passed from the InferType.  In $ce$ the
%%   unspecified widths are recorded by zero width .  The statements
%%   are assumed to contain only legal connections, which means the two
%%   sides of the connection have equivalent type (mentioned?) and the
%%   flow is from a smaller width to a lager one.
%% \item[The output] of InferWidth give a new component environment $ce'$
%%   where all unspecified types have been infered.  Still it is possible
%%   that after this pass, signals are infered to zero width.  The overall
%%   compilation will use other passes to deal with zero width, which is
%%   not concerned here.
%% \end{description}

%% To state the correctness of InferWidth that implement in Coq, we first
%% recall the semantics of width infering.  There are statements that
%% declare new component, or the statements that define connections. For
%% those components declared with explicit widths, we state that the
%% widths will be kept unchange after the pass.  The implicit width
%% declarations will be collected and await to be assigned to a proper width.
%% The connections give hints to the infering process. For example,
%% \begin{lstlisting}
%%   wire c: SInt
%%   c <= shl(SInt("b-100101100101"), 1)
%% \end{lstlisting}
%% the width of wire $c$ is implicit defined. The connection from an
%% shift-left expression to $c$ will provide the protential width for $c$.
%% According to the rules, the width of c now is 13, because the
%% shift-left expression is of type \lstinline!SInt<13>!. Before applying last
%% connect semantics, the module may contain more than one connections to
%% one target. For example,
%% \begin{lstlisting}
%%   wire c: SInt
%%   wire d: SInt<14>
%%   c <= shl(SInt("b-100101100101"), 1)
%%   c <= d
%% \end{lstlisting}
%% Now, wire $c$ has the width as same as $d$, which is 14. The latter connection
%% from d gives a larger width to the destination $c$.

%The bit widths of some components may still be missing after applying {\tt infer\_type} to all the statements in $ss$.
We formalize the ``InferWidths'' pass by defining a Coq function {\tt infer\_width}, which iterates over the statements in the sequence $ss$ and updates
the component environment $ce$, so that the bit widths of all the components with undefined bit widths can be determined.
%Note that after applying {\tt infer\_type} to all the statements in $ss$, if the bit width of a component is nonzero, that is, a bit width is available in its declaration statement, then its bit width will \emph{not} be changed by  {\tt infer\_width}.
The rationale behind the {\tt infer\_width} function is stated as follows: After applying the {\tt infer\_type} function to all the statements in the sequence $ss$, for each component $x$,
\begin{itemize}
\item if the bit width of the component $x$ is nonzero, that is, a bit width is available in its declaration statement, then
the function {\tt infer\_width} will \emph{not} change the bit width of the component $x$,
%
\item otherwise, the function  {\tt infer\_width} chooses the maximum width of the expressions that are connected to the component $x$.
%\item The last connection semantics is not considered at this stage,
%  all connect statements are taken into account.
\end{itemize}
Note that the last-connected semantics is ignored here. The behavior of the function {\tt infer\_width} on a connect statement is defined by the following rule,
%
  \begin{gather*}
    \inferiii[]
    {\seqq{ce \vdash {\tt infer\_width }(x<=e) \rightsquigarrow ce'}}
    {\seqq{ce(x) =  (t_x, \kappa)}}
    {\seqq{ce \vdash t_e = {\tt type\_of\_expr }(e)}}
    {\seqq{ce' = ce \uplus \{x : ({\tt max\_width }(t_x, t_e), \kappa)\}}}.
    \end{gather*}
%Intuitively, {\tt infer\_width} infers for each component the maximum width of the expressions that are connected to it. Note that the last-connected semantics is ignored here.
Moreover, the function {\tt infer\_width} uses another Coq function {\tt max\_width}, which, given two aggregate types, compares for each field their widths in the aggregate types, and chooses the larger width.
Finally,  the function {\tt infer\_width} processes the partial connect statements similarly.
%Note that {\tt infer\_width} does not check the equivalence of the bit widths of the two sides of connect statements.

%%%%%%%%%%%%%%%
%\hide{
%It is allowed to leave the width of a component implicit in HiFIRRTL programs.
%Such unspecified widths are inferred by the Infer Widths pass.
%The general rules are,
%\begin{itemize}
%\item The only allowed width for an explicitly defined width is itself.
%\item An implicitly defined width will be set to the largest width that can
%  be found through connect statements.
%\item The last connection semantics is not considered at this stage,
%  all connect statements are taken into account.
%\end{itemize}
%
%Following the rules, we first introduce an additional map $wm$, which
%maps from variables to types. Then we scan through the statements in $ss$ to
%find components with implicit width and create entries for
%them in $wm$. For a statement that declares a component $v$,
%the {\tt infer\_width} function will update $wm$ by
%%
%  \begin{gather}
%    \inferii[]
%    {\seqq{ce \vdash {\tt infer\_width} (declare\;v, wm) \rightsquigarrow wm'}}
%    {\seqq{ce \vdash {\tt width}(t_v) = 0}} {\seqq{wm':= wm \uplus \{v : t_v\}}}.
%    \end{gather}
%
%The width is inferred according to connect statements.
%Since the last connect semantics is not considered, each
%connection to the element provides a potential width. A function
%named {\tt max\_width} compares the potential widths and outputs
%the maximum one. For a statement which connects an expression $e$
%to a component $v$, the {\tt infer\_width} function will update $wm$ by
%
%After applying {\tt infer\_width} to all the statements in $ss$,
%$wm$ contains the inferred width for every implicitly defined
%component. By merging the information from $wm$ for every entry into $ce$,
%the component environment $ce$ contains the
%result of the inferred width.
%
%Instead of applying {\tt infer\_width} directly on $ce$, the map
%$wm$ contains only type information of the implicit ones that can be
%manipulated more efficiently.
%}
%%%%%%%%%%%%%%%

%% \begin{description}
%% \item[The input] of InferWidth takes the statments and the component
%%   environment $ce$ that passed from the InferType.  In $ce$ the
%%   unspecified widths are recorded by zero width .  The statements
%%   are assumed to contain only legal connections, which means the two
%%   sides of the connection have equivalent type (mentioned?) and the
%%   flow is from a smaller width to a lager one.
%% \item[The output] of InferWidth give a new component environment $ce'$
%%   where all unspecified types have been infered.  Still it is possible
%%   that after this pass, signals are infered to zero width.  The overall
%%   compilation will use other passes to deal with zero width, which is
%%   not concerned here.
%% \end{description}

%% To state the correctness of InferWidth that implement in Coq, we first
%% recall the semantics of width infering.  There are statements that
%% declare new component, or the statements that define connections. For
%% those components declared with explicit widths, we state that the
%% widths will be kept unchange after the pass.  The implicit width
%% declarations will be collected and await to be assigned to a proper width.
%% The connections give hints to the infering process. For example,
%% \begin{lstlisting}
%%   wire c: SInt
%%   c <= shl(SInt("b-100101100101"), 1)
%% \end{lstlisting}
%% the width of wire $c$ is implicit defined. The connection from an
%% shift-left expression to $c$ will provide the protential width for $c$.
%% According to the rules, the width of c now is 13, because the
%% shift-left expression is of type \lstinline!SInt<13>!. Before applying last
%% connect semantics, the module may contain more than one connections to
%% one target. For example,
%% \begin{lstlisting}
%%   wire c: SInt
%%   wire d: SInt<14>
%%   c <= shl(SInt("b-100101100101"), 1)
%%   c <= d
%% \end{lstlisting}
%% Now, wire $c$ has the width as same as $d$, which is 14. The latter connection
%% from d gives a larger width to the destination $c$.

\subsubsection{Formalization of InferResets}
We formalize the ``InferResets'' pass by defining a Coq function {\tt infer\_reset}, which iterates over the statements in
the sequence $ss$ and updates the component environment $ce$, so that the concrete reset types of all abstract reset signals can be inferred.
%Note that after applying {\tt infer\_type} to all the statements in $ss$, if the bit width of a component is nonzero, that is, a bit width is available in its declaration statement, then its bit width will \emph{not} be changed by  {\tt infer\_width}.
The rationale behind the function {\tt infer\_reset} is stated as follows: After applying the function {\tt infer\_type} to all the statements in the sequence $ss$, for each component $x$ of type \prettyll{Reset},
\begin{itemize}
\item if only components of  type \prettyll{AsyncReset} are connected to  the component $x$, then the function {\tt infer\_reset} updates the type of the component $x$ to \prettyll{AsyncReset},
%
\item if $x$ \prettyll{is invalid} is a statement in the sequence $ss$ or only components of type \prettyll{UInt<1>} are connected to the component $x$, then the function
{\tt infer\_reset} updates the type of the component $x$ to the type \prettyll{UInt<1>},
%
\item if two components of type \prettyll{AsyncReset} and \prettyll{UInt<1>} respectively are both connected to the component $x$, then the function  {\tt infer\_reset} reports an error.
\end{itemize}
The behavior of the  {\tt infer\_reset} function on a connect statement is defined by the following two rules,
 % \begin{gather*}
%    \label{rule-infer-reset-pass}
%    \inferiv[]
%    {\seqq{ce \vdash {\tt infer\_reset} (x<=e, ce) \rightsquigarrow ce'}}
%    {\seqq{ce(x)= (t_x, \kappa)}} {\seqq{ce\vdash t_e={\tt type\_of\_expr}(e)}}
%    {\seqq{{\tt \neg conflict} (t_x, t_e)}}
%    {\seqq{ce'= ce \uplus \{x : (t_e, \kappa)\}}}
%    \end{gather*}
and
  \begin{gather*}
    \inferiii[]
    {\seqq{ce \vdash {\tt infer\_reset} (x<=e, ce) \rightsquigarrow \bot}}
    {\seqq{ce(x) = (t_x, \kappa)}} {\seqq{ce\vdash t_e={\tt type\_of\_expr}(e)}}
    {\seqq{{\tt conflict} (t_x, t_e)}}.
%    {\seqq{ce'= \bot}}.
    \end{gather*}
Note that the function {\tt infer\_reset} uses the predicate {\tt conflict} to decide whether two reset types are conflicting, that is, one is  \prettyll{AsyncReset} and the other is \prettyll{UInt<1>}.
Moreover, here we take \prettyll{is invalid} statements as connect statements connected to a synchronous reset component (i.e. of type \prettyll{UInt<1>}).
%If there is no error reported, the inferred reset types stored in the temporary map $rm$ can be
%merged into the component environment $ce$ using the {\tt map2} function.
%The resulting $ce$ contains the inferred concrete reset types.

%%% the original texts written by xiaomu
%%% the original texts written by xiaomu
\hide{
In HiFIRRTL, the type \prettyll{Reset} can be used to indicate an
abstract reset signal. The Infer Reset pass is responsible for
inferring the concrete type (asynchronous or synchronous) of the abstract one based on the
connection context. The general rules are,

Similar to the Infer Widths pass, a temporary map $rm$ is introduced to
handle abstract resets exclusively. It maps from components to reset
types. The map $rm$ is first initialized with all components of type \prettyll{Reset}
by \eqref{rule-infer-reset-init}.
%% scanning through the statements in $ss$
%% for abstract reset types to be added to $rm$.
%% \david{I am not sure I understand the previous sentence:
%% is $rm$ filled based on $ss$ or based on $ce$?
%% Should it perhaps be:
%% $rm$ is first initialized with all (sub)components of type \prettyll{Reset} by rule \eqref{rule-infer-reset-init},
%% and then the pass scans through the statements in $ss$
%% to decide whether each one of them should become asynchronous or synchronous by rule \eqref{rule-infer-reset-pass}.}

  \begin{gather}
  \label{rule-infer-reset-init}
    \inferii[]
    {\seqq{ce \vdash {\tt infer\_reset} (declare\; v, rm) \rightsquigarrow rm'}}
    {\seqq{ce \vdash t_v : \mbox{Reset}}} {\seqq{rm':= rm \uplus \{v : t_v\}}}.
    \end{gather}
It then continues on the connect
statements to find the inference clues,
checking that the type of the driving signal and the type found in the $rm$ do not conflict.
If they are the same concrete type or at least one of them is abstract,
the inferring continues without error.
  \begin{gather}
    \label{rule-infer-reset-pass}
    \inferiv[]
    {\seqq{ce \vdash {\tt infer\_reset} (v<=e, rm) \rightsquigarrow rm'}}
    {\seqq{rm \vdash v:t_v}} {\seqq{ce\vdash t_e:={\tt type\_of\_expr}(e)}}
    {\seqq{{\tt \neg conflict} (t_v,t_e)}}
    {\seqq{rm':= rm \uplus \{v : t_e\}}}.
    \end{gather}
Otherwise, the {\tt infer\_reset} reports an error in the map $rm$ and stops.
  \begin{gather}
    \inferiv[]
    {\seqq{ce \vdash {\tt infer\_reset} (v<=e, rm) \rightsquigarrow rm'}}
    {\seqq{rm \vdash v:t_v}} {\seqq{ce\vdash t_e:={\tt type\_of\_expr}(e)}}
    {\seqq{{\tt conflict} (t_v,t_e)}}
    {\seqq{rm':= \bot}}.
    \end{gather}
Here, we consider the case of \prettyll{is invalid} statements as same as a connect statement from a synchronous signal.
If there is no error reported, the inferred reset types stored in the temporary map $rm$ can be
merged into the component environment $ce$ using the {\tt map2} function.
The resulting $ce$ contains the inferred concrete reset types.
}
%%% the original texts written by xiaomu
%%% the original texts written by xiaomu

\subsubsection{Formalization of ExpandConnects}
The ``ExpandConnects'' pass expands the connections between aggregate
types into those between ground types. It also flattens the elements
in aggregate types to their corresponding ground types in declaration
statements. To expand the objects of aggregate types in a Firrtl
module, the process involves two phases. The first phase is to record
the types along with orientations and flipping information of the
objects, which helps determine the actual connection directions. For
example, if there is an input port of aggregate types with a field
marked as ``flip'', the ``flip'' mark will recursively apply to all the
atomic elements in that field. Since it is an input port, the
orientation is record as ``source''. The second phase is to expand the statements
sequentially. For declaration statements, ground-typed elements are
generated corresponding to the original aggregate types. For
connection statements, the connection directions are determined by the
orientations and flipping information from the first phase. For
subfields, there are two kinds of connections allowed for HiFirrtl
programs. One is to set connections between elements of equivalent
aggregate types. In such bundled connections, individual subfield
connections can be of different directions. The other kind is to
connect subfields separately, in which case the direction should
comply with the orientation of the subfield elements.


At the beginning of a Firrtl module, ports are declared first. When
ports are defined with flip sub-fields, the expansion process will
change those fields of input ports into output ports, and vice
versa. As a result, their orientations are changed between ``source''
and ``sink''.
 \begin{gather*}
   \label{expand-input}
   \infer[]
         {\seqq{{\tt expand\_connect} (\mbox{\prettyll{input v}}) = ss}}
         {\begin{array}{c}
             \seqq{{\tt expand\_ref \;v} = [v_0, v_1,\dots v_n]}\\
             \seqq{ s s= {\tt map}\; (fun\; v_x \Rightarrow if\;{\tt is\_flipped}(v_x)\;then\; \mbox{\prettyll{output}} {\tt v_x}\;else\; \mbox{\prettyll{input}} {\tt v_x})\, [v_0, v_1,\dots v_n]}
           \end{array}
         }
 \end{gather*}

Elements such as wires and registers have duplex
orientation. Regardless of the flip sub-fields in the original
aggregate types, the resulting orientations will remain duplex. To
expand registers in aggregate types, we also need to expand the reset
values of the registers. Expanding wires is relatively simple and is
omitted here.
\begin{gather*}
  \label{expand-reg}
  \infer[]
        {\seqq{{\tt expand\_connect} (\mbox{\prettyll{reg v : t clk with : (reset => rst, rv)}}) = ss}}
        {\begin{array}{c}
            \seqq{{\tt expand\_ref \;}v = [v_0,\dots v_n]}\\
            \seqq{{\tt expand\_expr \;}rv = [rv_0,\dots rv_n]}\\
            \seqq{ s s= {\tt map}\; (fun\; v_x\; rv_x \Rightarrow \mbox{\prettyll{reg v_x: t clk with:(reset=> rst,rv_x)}})}
            %% \rightline{
              \seqq{[v_0,\dots v_n], [rv_0,\dots rv_n]}
          \end{array}
        }
\end{gather*}

Nodes are also ``source'' elements. They are only allowed to be declared
with passive aggregate types, which means no flip sub-fields. During
the expansion of node declarations, the passive type is checked;
otherwise, the result will be an empty sequence.
\begin{gather*}
  \label{expand-node}
  \infer[]
        {\seqq{{\tt expand\_connect} (\mbox{\prettyll{node v e}}) = ss}}
        {\begin{array}{c}
            \seqq{{\tt expand\_ref}\; v = [v_0, v_1,\dots v_n]}\\
            \seqq{{\tt expand\_expr}\; e = [e_0, e_1,\dots e_n]}\\
            \seqq{ s s= if\;{\tt is\_passive}(v_x)\; then\; {\tt map}\; (fun\; v_x\;e_x \Rightarrow \; \mbox{\prettyll{node}} {\tt v_x}\;{\tt e_x})\,[v_0, v_1,\dots v_n]\,[e_0, e_1,\dots e_n] else\; [::]}
          \end{array}
        } 
\end{gather*}

To expand \prettyll{is invalid} statements, we
select only the ``sink'' and ``duplex'' elements to produce the sequence of
invalidations.
In HiFirrtl, it is allowed to have ``source'' sub-elements in a bundle be
invalidated. The result of expansion will omit them.
\begin{gather*}
  \label{expand-invalid}
  \infer[]
        {\seqq{{\tt expand\_connect} (\mbox{\prettyll{v is invalid}}) = ss}}
        {\begin{array}{c}
            \seqq{{\tt expand\_ref \;}v = [v_0, v_1,\dots v_n]}\\
            \seqq{vs = {\tt filter}\; (fun\; v_x \Rightarrow orient (v_x) \neq source) [v_0, v_1,\dots v_n]}\\
            \seqq{ s s = {\tt map}\; (fun\; v_x \Rightarrow {\tt v_x}\;\mbox{\prettyll{is invalid}})\, vs}
          \end{array}
        } 
\end{gather*}

Connections between bundled elements contain a set of connections for
individual sub-fields. According to the indicated ``flip'' for a sub-field, the
connection direction may differ from the original.
\begin{gather*}
   \label{expand-invalid}
   \infer[]
         {\seqq{{\tt expand\_connect} (\mbox{\prettyll{l <= r}}) = ss}}
         {\begin{array}{c}
             \seqq{{\tt expand\_ref }\; l = [l_0,\dots l_n]}\\
             \seqq{{\tt expand\_expr}\; r = [r_0,\dots r_n]}\\
             \seqq{ s s = {\tt map}\; (fun\; l_x\;r_x \Rightarrow if\;{\tt is\_flipped}(l_x)\; then\;\mbox{\prettyll{r_x <= l_x}}\; else\mbox{\prettyll{l_x <= r_x}})\,[l_0,\dots l_n]\,[r_0,\dots r_n]}
           \end{array}
         }
\end{gather*}




\hide{
The ``ExpandConnects'' pass expands the connections between aggregate types into those between ground types.
We define two Coq functions {\tt expand\_expr} and {\tt expand\_cnct} that expand the expressions and the connect statements, respectively.
The rationale behind the functions {\tt expand\_expr} and {\tt expand\_cnct} is to break aggregate types in expressions and connect statements recursively into their fields.
Note that the function {\tt expand\_cnct} expands both connect and partial connect statements. The functions {\tt expand\_expr} and {\tt expand\_cnct} modify the sequence $ss$ of statements,
instead of component environment $ce$.
%\begin{itemize}
%\item connections between ground types can not be expanded,
%
%\item connections between aggregate types will be expanded recursively.
%\end{itemize}

%In HiFIRRTL, connections can be made between aggregate types.  These
%connections can be either full or partial.  The purpose of this pass
%is to break down connections between aggregate types and produce
%individual ones of corresponding ground types. The general rules are,
%\begin{itemize}
%\item Connections between ground types can not be expanded.
%\item Connections between aggregate types will be expanded recursively.
%\end{itemize}

The function {\tt expand\_expr} expands an expression $e$ into a list of expressions of ground types recursively by utilizing {\tt type\_of\_expr}$(e)$.
Note that only expressions using the operators of multiplex \prettyll{mux} and
 \prettyll{validif}, or references can have aggregate types, while the expressions formed by using binary operators e.g. $add$ do not.
The function {\tt expand\_expr} expands the \prettyll{mux} and
 \prettyll{validif} expressions by the following two rules,
  \begin{gather*}
  \label{expand-mux}
    \inferiii[]
    {\seqq{ce\vdash{\tt expand\_expr} (\mbox{\prettyll{mux}}(e_c, e_t, e_f)) \rightsquigarrow els}}
    {\seqq{els = \mbox{\tt ls\_mux}(e_c, els_t, els_f)}}
    {\seqq{els_t = {\tt expand\_expr}(e_t)}}
    {\seqq{els_f = {\tt expand\_expr}(e_f)}}
    \end{gather*}
    and
  \begin{gather*}
  \label{expand-validif}
    \inferii[]
    {\seqq{ce\vdash{\tt expand\_expr} (\mbox{\prettyll{validif}}(e_c, e)) \rightsquigarrow els}}
    {\seqq{els  = \mbox{\tt ls\_validif}(e_c, els_t)}}
    {\seqq{els_t = {\tt expand\_expr}(e)}}.
    \end{gather*}
The first rule %~(\ref{expand-mux})
uses a Coq function {\tt ls\_mux} that combines two sequences of expressions $els_t$ and $els_f$ with the same condition $e_c$, resulting into a list of ground-type expressions
$$\mbox{\prettyll{mux}}(e_c,els_{t, 1}, els_{f, 1}), \cdots, \mbox{\prettyll{mux}}(e_c, els_{t, m}, els_{f, m}),$$
where  $els_{t,i}$ (resp. $els_{f, i}$) denotes the $i$-th expression in the sequence of expressions $els_t$ (resp. $els_f$).
The function {\tt ls\_validif} in the second rule %~(\ref{expand-validif})
is defined similarly.
% to the $i$-th element of $els_t$ and $els_f$ .
%
%It takes the conditional expression $e_c$ and the true/false branch
%expressions $e_t$ and $e_f$ from the original \prettyll{mux} expression, and then produces a list $es_{mux}$
%of \prettyll{mux} that combines the lists of expanded $e_t$ and $e_f$
%with the same condition $e_c$.
%Similarly, expanding a \prettyll{validif} expression produces a list of \prettyll{validif} expressions in $es_{validif}$,
%
Finally, {\tt expand\_expr}($r$) expands a reference $r$ straightforwardly and creates a list of expressions of ground types according to {\tt type\_of\_expr$(r)$}.
%Expanding a references is relatively straightforward.
%It create a list of references to the sub-elements according to the
%type of the component.

Equipped with the {\tt expand\_expr} function, the function {\tt expand\_cnct} works as follows: Given a connect statement $(x <= e)$, it applies
the function {\tt expand\_expr} to the component $x$ and the expression $e$ respectively, producing two sequences of expressions of ground types, then constructs a sequence of connect statements by combining the expressions in the two sequences. The behavior of the {\tt expand\_cnct} function on a connect statement is defined by the following rule,
%After expanding the list of references on the left-hand side and the
%list of expressions on the right-hand side, new connect statements can
%be generated pairwise from the two by
  \begin{gather*}
  \label{expand-cnct}
    \inferiii[]
    {\seqq{ce\vdash{\tt expand\_cnct} (x <= e) \rightsquigarrow ss'}}
    {\seqq{ ss' = {\tt ls\_cnct}(els_l, els_r)}}
    {\seqq{els_l = {\tt expand\_expr}(x)}}
    {\seqq{els_r = {\tt expand\_expr}(e)}}.
    \end{gather*}
The function {\tt expand\_cnct} uses a Coq function {\tt ls\_cnct} that combines the two sequences $els_l$ and $els_r$ of expressions into a sequence $ss'$ of connect statements between ground types, that is,
\begin{center}
the statements $els_{l, 1} <= els_{r, 1}, \cdots, els_{l, m} <= els_{r, m}$,
\end{center}
 where $m$ is the length of the two sequences, and $els_{l, i}$ (resp. $els_{r, i}$) denotes the $i$-th expression of the sequence $els_l$ (resp. $els_r$).

The function {\tt expand\_cnct} handles the \prettyll{is invalid} statements in the same way as connect statements. It produces from a \prettyll{is invalid} statement
 a sequence of \prettyll{is invalid} statements for its fields.
 }

%%%%%%%%%%%%%%%%%%%%%%%%%%%%%%
%%%%%%%%%%%%%%%%%%%%%%%%%%%%%%
%\hide{
%The pass focuses on connect statements in $ss$.
%The left-hand side
%of a connection is a reference to a component, while the right-hand side
%is an expression. In order to obtain expanded
%connections, the expression needs to be recursively expanded.
%Only expressions of multiplexers \prettyll{mux},
%\prettyll{validif} and
%references \prettyll{r} can be of aggregate types. The other
%expressions, such as binary operations, only operate and result in
%ground types. Therefore they do not require expanding.
%Expanding a \prettyll{mux} expression is done by,
%  \begin{gather}
%  \label{expand-mux}
%    \infer[]
%    {\seqq{ce\vdash{\tt expand\_expr} (\mbox{mux}(e_c,e_t,e_f)) \rightsquigarrow es_{mux}}}
%    {\seqq{es_{mux}:=[\mbox{mux}(e_c,t,f)\,|\,t \leftarrow {\tt expand\_expr}(e_t),\,f \leftarrow {\tt expand\_expr}(e_f)]}}.
%    \end{gather}
%It takes the conditional expression $e_c$ and the true/false branch
%expressions $e_t$ and $e_f$ from the original \prettyll{mux} expression, and then produces a list $es_{mux}$
%of \prettyll{mux} that combines the lists of expanded $e_t$ and $e_f$
%with the same condition $e_c$.
%Similarly, expanding a \prettyll{validif} expression produces a list of \prettyll{validif} expressions in $es_{validif}$,
%  \begin{gather}
%  \label{expand-validif}
%    \infer[]
%    {\seqq{ce\vdash{\tt expand\_expr} (\mbox{validif}(e_c,e)) \rightsquigarrow es_{validif}}}
%    {\seqq{es_{validif}:=[\mbox{validif}(e_c,t)\,|\,t \leftarrow {\tt expand\_expr}(e)]}}.
%    \end{gather}
%Expanding a references is relatively straightforward.
%It create a list of references to the sub-elements according to the
%type of the component.
%After producing the list of references on the left-hand side and the
%list of expressions on the right-hand side, new connect statements can
%be generated pairwise from the two by
%  \begin{gather}
%  \label{expand-cnct}
%    \infer[]
%    {\seqq{ce\vdash{\tt expand\_cnct} (v <= e) \rightsquigarrow ss_{cnct}}}
%    {\seqq{ ss_{cnct}:=[lhs<=rhs\,|\,lhs \leftarrow {\tt expand\_ref}(v),\,rhs \leftarrow {\tt expand\_expr}(e)]}}.
%    \end{gather}
%
%The \prettyll{is invalid} statement is treat as the same as the connect statement,
%which connects a component with an ``unknown'' signal at this point.
%After processing an \prettyll{is invalid} statement for a component,
%we obtain a list of \prettyll{is invalid} statements for its sub-elements.
%}


%%%%%%%% remove resets removed
%%%%%%%% remove resets removed
%\hide{
%\subsubsection{``Remove Resets''}
%The Remove Resets pass runs after the other passes, and removes
%redundant reset that are synchronous. This is a optimization step that
%can potentially improve the performance by removing unnecessary ones.
%A synchronous reset is considered redundant if it is only used to
%initialize a register.
%}
%%%%%%%% remove resets removed
%%%%%%%% remove resets removed
%% For example,
%% \begin{lstlisting}
%% module MyModule :
%%   input myinput: {a: UInt<1> [1]}
%%   output myoutput: {a: UInt<1> [2], b: UInt<1>}
%%   myoutput <- myinput
%% \end{lstlisting}
%% the connection between myinput and myoutput is a partial connection of
%% aggregate types. It can then be expand through ExpandConnects pass and get,
%% \begin{lstlisting}
%% module MyModule :
%%   input myinput_a: UInt<1>
%%   output myoutput_a_0: UInt<1>
%%   output myoutput_a_1: UInt<1>
%%   output myoutput_b: UInt<1>
%%   myoutput_a_0 <= myinput_a_0
%% \end{lstlisting}

%% \begin{description}
%% \item[The input] of ExpandConnects consists of the sequence of
%%   statments representing the original program, the result component
%%   enviroment $ce3$ from InferWidths pass. The $ce3$ collects the
%%   component identifies to their types.
%% \item[The output] of ExpandConnects contains a program in which all
%%   connections between aggregate types are replaced by connections of
%%   their sub elements and a corresponding $ce4$. In $ce4$, new
%%   identifiers added and their types are consistent to the flattened
%%   aggregate type. As the aboved example, ouput port \lstinline!myport!
%%   is replaced by 3 output ports of ground types.
%% \end{description}

\subsubsection{Formalization of ExpandWhens}
The ``ExpandWhens'' pass ensures that every component in a module is connected to at most once.
It reformulates the conditions in the \prettyll{when} statements by multiplex expressions and removes the \prettyll{when} statements. Moreover, for each component, it removes the non-last (redundant) connect statements for the component.
We define two functions {\tt expand\_branch} and {\tt combine\_branches} to formalize the ``ExpandWhens'' pass.
The function {\tt expand\_branch} scans through the statement sequence $ss$, copies the declaration statements, and updates the component environment $cm$ when processing connect, \prettyll{is invalid}, and \prettyll{when} statements.  The function {\tt combine\_branches} is used to combine the connect mappings for branches of \prettyll{when} statements.
%This is done by scanning through the statement sequence of the module:
%declaration statements are copied without change,
%and connect and \prettyll{is invalid} statements update a connect map
%(so redundant connects are overwritten in the map).
The  {\tt expand\_branch} function is defined by the following four rules, % (\ref{rule-expandwhen-declare})-(\ref{rule-expandwhen-when}),
where $\mathit{declare}\ x$ represents a declaration statement for $x$, e.g. \prettyll{wire x: t}.
\begin{gather*}
    \label{rule-expandwhen-declare}
    \infer[]
    {\seqq{{\tt expand\_branch}(ss ; \mathit{declare}\ x, cm_0) = (ss_\mathrm{dcl} ; \mathit{declare}\ x, cm[\bot/x])}}
    {\seqq{{\tt expand\_branch}(ss, cm_0) = (ss_\mathrm{dcl}, cm)}}
\\[1ex]
    \label{rule-expandwhen-connect}
    \infer[]
    {\seqq{{\tt expand\_branch}(ss ; x <= e, cm_0) = (ss_\mathrm{dcl}, cm [e/x])}}
    {\seqq{{\tt expand\_branch}(ss, cm_0) = (ss_\mathrm{dcl}, cm)}}
\\[1ex]
    \label{rule-expandwhen-invalid}
    \infer[]
    {\seqq{{\tt expand\_branch}(ss ; v \text{ \prettyll{is invalid}}, cm_0) = (ss_\mathrm{dcl}, cm [{\tt invalid}/v])}}
    {\seqq{{\tt expand\_branch}(ss, cm_0) = (ss_\mathrm{dcl}, cm)}}
\end{gather*}
%
\begin{gather*}
    \label{rule-expandwhen-when}
    \inferii[]
    {\seqq{{\tt expand\_branch}(ss ; \text{\prettyll{when c then} } ss^\mathsf{T} \text{ \prettyll{else} } ss^\mathsf{F}, cm_0) =
        (ss_\mathrm{dcl} ; ss_\mathrm{dcl}^\mathsf{T} ; ss_\mathrm{dcl}^\mathsf{F}, {\tt combine\_branches}(cm^\mathsf{T}, cm^\mathsf{F}))}}
    {\seqq{{\tt expand\_branch}(ss, cm_0) = (ss_\mathrm{dcl}, cm)}}
    {\begin{gathered}[b] \seqq{{\tt expand\_branch}(ss^\mathsf{T}, cm) = (ss_\mathrm{dcl}^\mathsf{T}, cm^\mathsf{T})} \\
                      \seqq{{\tt expand\_branch}(ss^\mathsf{F}, cm) = (ss_\mathrm{dcl}^\mathsf{F}, cm^\mathsf{F})} \end{gathered}}
\end{gather*}
Note that for a statement \prettyll{when c then} $ss^\mathsf{T}$ \prettyll{else} $ss^\mathsf{F}$,
 the function {\tt expand\_branch}  first processes the statements preceding this statement and obtains a connect mapping $cm$, then uses $cm$ to process $ss^\mathsf{T}$ and $ss^\mathsf{F}$ respectively, resulting into the connect mappings $cm^\mathsf{T}$ and $cm^\mathsf{F}$ respectively. Finally, $cm^\mathsf{T}$ and $cm^\mathsf{F}$ are combined into one mapping by {\tt combine\_branches}.

%For \prettyll{when} statements,
%the true and false branches are similarly split
%and the two resulting connect maps are then combined into one,
%using multiplexer or \prettyll{validif} expressions where necessary.

%In contrast to the operational LoFIRRTL semantics defined in Subsection~\ref{sec-lofirrtl-sem},
%the connect statements do not need to be ordered topologically.
%Therefore, we cannot evaluate expression $e$ in rule~\ref{rule-expandwhen-connect} immediately
%but copy it as a syntactical object into the connect map.
%This also implies that we cannot use the (register) state of a configuration in these rules.


The function {\tt combine\_branches} combines two connect mappings into one by the following fix rules %(\ref{rule-expandwhen-combine-1})-(\ref{rule-expandwhen-combine-6}),
where an expression is constructed for each component, denoting the fact that the expression is connected to it according to the last-connect semantics.
If both connect mappings assign the same value to a component, then no multiplexer or \prettyll{validif} is needed in the combined expression.
\begin{gather*}
 \label{rule-expandwhen-combine-1}
    \infer[]
    {{\tt combine\_branches}(cm^\mathsf{T},cm^\mathsf{F})(x) = e}
    {cm^\mathsf{T}(x) = cm^\mathsf{F}(x) = e}
\\[1ex]
 \label{rule-expandwhen-combine-2}
    \infer[]
    {{\tt combine\_branches}(cm^\mathsf{T},cm^\mathsf{F})(x) = {\tt invalid}}
    {cm^\mathsf{T}(x) = cm^\mathsf{F}(x) = {\tt invalid}}
\end{gather*}
If the two connect mappings assign different values to a component, then a multiplexer or \prettyll{validif} expression is obtained.
\begin{gather*}
 \label{rule-expandwhen-combine-3}
    \inferiii[]
    {{\tt combine\_branches}(cm^\mathsf{T},cm^\mathsf{F})(x) = \text{\prettyll{mux}}(\text{\prettyll{c}},e^\mathsf{T},e^\mathsf{F})}
    {cm^\mathsf{T}(x) = e^\mathsf{T}}
    {cm^\mathsf{F}(x) = e^\mathsf{F}}
    {e^\mathsf{T} \not= e^\mathsf{F}}
\\[1ex]
 \label{rule-expandwhen-combine-4}
    \inferii[]
    {{\tt combine\_branches}(cm^\mathsf{T},cm^\mathsf{F})(x) = \text{\prettyll{validif}}(c,e^\mathsf{T})}
    {cm^\mathsf{T}(x) = e^\mathsf{T}}
    {cm^\mathsf{F}(x) = {\tt invalid}}
\end{gather*}
If only one branch declared the component, ignore the condition;
if at least one branch did declare the component but erroneously did not connect it,
register that it is still not (always) connected:
\begin{gather*}
 \label{rule-expandwhen-combine-5}
    \inferii[]
    {{\tt combine\_branches}(cm^\mathsf{T},cm^\mathsf{F})(x) = e^\mathsf{T}}
    {cm^\mathsf{T}(x) = e^\mathsf{T}}
    {cm^\mathsf{F}(x) \text{ is undefined}}
\\[1ex]
 \label{rule-expandwhen-combine-6}
    \infer[]
    {{\tt combine\_branches}(cm^\mathsf{T},cm^\mathsf{F})(x) =\bot}
    {cm^\mathsf{F}(x) = \bot}
\end{gather*}
For the last three rules, a symmetric variant exists where the role of the true and false branches are swapped.

Moreover, to initialize, the empty statement sequence does not change a connect map, specified by the following rule:
{\begin{gather*}
    \infer{\seqq{{\tt expand\_branch}(\emptyset, cm_0) = (\emptyset, cm_0)}}
    {} % no premise
\end{gather*}}

Finally, the statement sequence is updated by concatenating the statement sequence generated by the function  {\tt expand\_branch} and the sequence of connect and \prettyll{is invalid} statements corresponding to the connect mapping generated by the function {\tt expand\_branch}.
%(possibly after topologically sorting them),
%similar to the example in Subsection~\ref{subsec-lowering-transformations}.
If in the end some component is still assigned the value $\bot$ by the connect mapping, an appropriate error message can be generated.

\subsubsection{Formalization of LowerTypes}
The ``LowerTypes'' pass flattens each component of an aggregate type into a collection of components of ground types.
%The pass replaces certain defined components with
%aggregate types by a new set of components defined with ground types.
We formalize this pass by a function {\tt lower\_types}.
The rationale behind the function {\tt lower\_types} is that except for the memories and module instances, each component of an aggregate type is replaced by
      a collection of new components that have ground types, one for each field, where \emph{paired identifiers} are introduced for  naming the new components.
%\begin{itemize}
%\item components defined with ground types do not need to be modified,
%
%\item Except for memories and module instances, components defined with aggregate types need to be replaced by
%      a group of coordinate components that have ground types.
%\end{itemize}

A paired identifier is of the form $(id_1, id_2)$, where $id_1$ is the identifier of a component of aggregate type, and $id_2$ is the name of a field in the aggregate type.
%the first field represents the base reference adapted from the original
%component identifier and the second field indicates the
%sub-element offset computed from the aggregate type.
%The original components defined with aggregate types are replaced by a new
%set of components defined with paired identifiers and of ground types.
The function {\tt lower\_types} processes the statements in the sequence $ss$ and updates the sequence  $ss$ and the component environment $ce$.
The behavior of the function {\tt lower\_types} on the declaration statements is defined by the rule below, where $declare\ x$ denotes the declaration statement for a component $x$, and $ce=(t, \kappa)$ ($t$ is the aggregate type and $\kappa$ is the kind, e.g. \prettyll{wire}), moreover, a function  {\tt flatten} is used to flatten $declare\ x$ into a sequence of declaration statements.  Specifically, let the fields in $t$ be $f_1, \cdots, f_n$, of ground types $t_1, \cdots, t_n$ respectively, then the {\tt flatten} function generates from $(x, (t, \kappa))$ a sequence of declaration statements $\kappa\ (x, f_1): t_1, \cdots, \kappa\ (x, f_n): t_n$.
%The process can be formalized as follow:
  {\begin{gather*}
    \inferiii[]
    {\seqq{ce\vdash{\tt lower\_types}(declare\;x) \rightsquigarrow (ss',ce')}}
    {\seqq{ce(x) = (t, \kappa)}}
%    {\seqq{{\tt is\_aggr}(t_v)}}
    {\seqq{ ss' ={\tt flatten}(x, (t, \kappa))}}
    {\seqq{ce'=ce\uplus \{(x, f_1): (t_{1},\kappa), \dots,(x, f_n):(t_{n}, \kappa)\}}}.
    \end{gather*}}
%The {\tt flatten} function generates new statements that define
%components in ground types with new identifiers, such as
%$(v,o_1)$. The type and component information contained in $tc_v$
%are used to determine the number of the certain new components.

Note that the ``LowerTypes'' pass should be applied after the ``ExpandConnects'' pass.
Therefore, before applying the function {\tt lower\_types} to the statements in the sequence $ss$, both sides of the connect statements are already in ground types.
Then when the {\tt lower\_types} function  processes a connect statement, the references to the component fields in this statement should be replaced by the corresponding paired identifiers. Thus in this case, all the three components of the circuit diagram, namely, the statement sequence, the component environment, the connect mapping, are all updated to reflect the replacements of component fields by paired identifiers.
%As a result, the statements $ss$ in the circuit diagram are replaced by the new $ss'$, and the components environment $ce$ are updated to reflect the new ground types and paired identifiers.

